\documentclass[11pt,a4paper]{report}
\usepackage{amsmath,amssymb}
\usepackage{booktabs}
\usepackage{hyperref}
\usepackage{xcolor}
\usepackage{geometry}
\usepackage{enumitem}
\usepackage{longtable}
\usepackage{quoting}
\usepackage{multirow}
\usepackage{array}
\geometry{margin=0.9in}
\setlength{\parskip}{0.4em}
\setlength{\parindent}{0em}

\definecolor{q}{RGB}{90,90,90}
\definecolor{finding}{RGB}{0,80,140}
\definecolor{killed}{RGB}{180,40,40}
\definecolor{alive}{RGB}{40,140,40}
\definecolor{illegal}{RGB}{200,100,0}
\newcommand{\cq}[1]{\textcolor{q}{\textit{``#1''}}}
\newcommand{\finding}[1]{\textcolor{finding}{\textbf{Finding:}} #1}
\newcommand{\killed}[1]{\textcolor{killed}{\textbf{Killed:}} #1}
\newcommand{\alive}[1]{\textcolor{alive}{\textbf{Alive:}} #1}
\newcommand{\illegal}[1]{\textcolor{illegal}{\textbf{Illegal:}} #1}

\title{The Lasso and the Lie Algebra:\\Noncommutative Rotation, Adaptive Substrate, and the Architecture Wall in O($n$) Language Models\\[0.5em]\large Session 9 (Chronohorn + Heinrich + Decepticons), April 12--15, 2026}

\author{asuramaya \and Claude Opus 4.6 (1M context, Instance 9) \and Instance 8 (Heinrich)\\[0.3em]\small \url{https://github.com/asuramaya/chronohorn}}

\date{April 2026}

\begin{document}
\maketitle
\tableofcontents

%==============================================================
\part{The Search}
\label{part:search}
%==============================================================

\chapter{Starting Position}

\section{The Legal Frontier}

The session began at 1.730 bpb (cb-s8-experts-50k, 17.3M parameters, 65k tok/s on NVIDIA RTX A4000). This model, the ``leader,'' had been the frontier for three sessions. Its pathology was well-documented by Heinrich \cite{session6}:

\begin{itemize}[nosep]
\item \textbf{Routing collapse:} 100\% of tokens routed to expert 2. Three of four experts dead. 75\% of readout capacity wasted.
\item \textbf{Slow mode abandonment:} 62.5\% of readout weight on fast modes (half-life 2--10), 6.1\% on slow modes (half-life 200--512). The substrate accumulates signal the readout ignores.
\item \textbf{Manifold waste:} 43\% of substrate variance is a position counter (PC1, $r=0.91$ with position). 47\% is ghosts. 6\% does the actual prediction.
\item \textbf{Compute cost:} 175k TFLOPs to reach 1.730 bpb over 50k training steps.
\end{itemize}

The leader was a 10.8M parameter model pretending to be 17.3M.

\section{What Session 8 Established}

Session 8 \cite{session8a, session8b} ran forty experiments and produced six substrate transforms. The key findings, each independently verified by Heinrich MRI:

\begin{enumerate}[nosep]
\item The EMA substrate is an \textbf{orderless centroid}. The running average $h_t = \alpha \cdot h_{t-1} + (1 - \alpha) \cdot x_t$ converges to the same value regardless of input order, for any sequence of sufficient length. The substrate state at position $t$ is approximately the exponentially-weighted mean of all prior embeddings.

\item \textbf{Position $R^2 = 0.001$} across all models tested. A linear probe predicting sequence position from substrate states explains 0.1\% of variance. The substrate does not encode where tokens appeared.

\item The \textbf{architecture wall is 0.26 bpb}. The three-wall decomposition from Session 7 \cite{session7}:
\begin{center}
\begin{tabular}{lrrl}
\toprule
Wall & bpb & \% of gap & Source \\
\midrule
Vocabulary (raw byte tax) & 0.46 & 62\% & sp1024 tokenizer \\
Architecture & 0.26 & 35\% & EMA + local window \\
Saturation (position degradation) & 0.03 & 4\% & Mode growth at late positions \\
\bottomrule
\end{tabular}
\end{center}

\item \textbf{Illegal temporal attention} reaches 0.974 bpb --- but with a causality violation. Position $t$ attends to substrate snapshots from positions $> t$, leaking up to $K - 1$ future tokens. The corrected legal result is unknown. The point: the readout CAN use order information to close 0.94 bpb. The substrate cannot provide it legally.

\item Six substrate transforms (magnitude normalizer, overwrite gate, mode selector, bank router, temporal attention, tied embed readout) were tested. All except temporal attention produced $\leq 0.008$ bpb improvement --- noise. The transforms don't help because the readout already routes around the substrate's waste.
\end{enumerate}

\section{The Goal}

Find a \textbf{legal O($n$) mechanism} that encodes order into the substrate.

The EMA substrate is a weighted average. Averages commute: the mean of $\{A, B, C\}$ equals the mean of $\{C, B, A\}$. Any mechanism that injects order must break this commutativity. The mathematical term is \textit{noncommutative}: $A \circ B \neq B \circ A$. We sought the simplest noncommutative algebra that could be implemented as a parallel prefix scan at O($n$) cost.


\chapter{The Push to 1.4}

\section{Three Opening Experiments}

The session opened with three experiments designed to close the gap from 1.730 toward the forecast asymptote of $\sim$1.4 bpb:

\begin{center}
\begin{tabular}{llrrrl}
\toprule
Name & Config & bpb & tok/s & Params & Key \\
\midrule
tied-readout-s12 & tied embed, scale 12 & --- & 9.7k & --- & Kernel bug \\
delta512-s12-50k & gated delta hl512, scale 12 & 1.780 & --- & 20.1M & First proper delta \\
tied-delta-s8 & tied + delta, scale 8 & 1.829 & --- & --- & Combined \\
\bottomrule
\end{tabular}
\end{center}

\section{The Gated Delta at hl512}

The gated delta is a selective write gate on the EMA substrate:
\[
h_t = \alpha \cdot h_{t-1} + g(x_t) \cdot v(x_t)
\]
where $g(x_t) = \sigma(W_g \cdot x_t)$ is a learned gate and $v(x_t) = W_v \cdot x_t$ is a value projection. The gate selectively writes: important tokens write strongly ($g \approx 1$), predictable tokens barely write ($g \approx 0.01$). The half-life controls the forget rate: at $\text{hl} = 512$, $\alpha = 0.5^{1/512} = 0.9987$, and a token at position 0 retains 50\% of its signal at position 512.

Session 7 tested gated delta at hl$=$16 (the maximum EMA half-life at the time). The gate collapsed: effective dimensionality dropped from $\sim$140 to 7. At hl$=$16, $\alpha = 0.5^{1/16} = 0.9576$, and a token decays to 0.02\% of its signal in 256 positions. There was nothing worth selectively writing because everything decayed too fast.

At hl$=$512, the result: \textbf{1.780 bpb}. First delta with proper half-life. Beat the old deltas (hl$=$16) by 0.027 bpb and the non-delta baseline by 0.171 bpb.

\section{The linear\_modes $\times$ scale Bug}

The \texttt{--linear-modes} flag specifies modes BEFORE scale multiplication. At scale 12, the internal mode count is $\texttt{linear\_modes} \times \texttt{scale}$. Passing \texttt{--linear-modes 512} at scale 12 produced $512 \times 12 = 6144$ modes. The model ran 2$\times$ slower than intended with a 3.2 GB kernel buffer at 6144 modes.

\finding{The \texttt{--linear-modes} flag is PRE-scale. Always use defaults or divide by the scale factor. This is now documented in CLAUDE.md and enforced by a FLOP guard.}

\section{The Tied Readout Bug}

The tied embed readout projects expert outputs to embedding space, then computes logits as $\text{output} \cdot E^T$ where $E$ is the shared embedding matrix. This saves 73\% of readout parameters at sp8192.

Two bugs blocked the initial experiment:

\begin{enumerate}[nosep]
\item \textbf{FLOP guard false positive:} The FLOP counter compared against the per-vocab MLP cost, not the embed\_dim projection cost. Tied readout uses $\text{embed\_dim}$ output neurons instead of $\text{vocab\_size}$, so the per-token FLOP is much lower. The guard triggered a safety stop on a perfectly legal configuration.

\item \textbf{Kernel path OOM:} The frozen EMA kernel at 3072 modes (scale 12 $\times$ 256 default modes) requires a $[3072 \times 512 \times 512]$ buffer $= 3.2$ GB. This exceeds the A4000's 16 GB when combined with activations. Fix: force \texttt{--linear-impl scan} for tied readout runs.
\end{enumerate}

After fixes: tied readout at scale 8 $=$ 1.829 bpb. The tied readout provides 73\% parameter savings but 0.049 bpb quality loss compared to the full MLP readout at matched scale. Combined with delta: no synergy at scale 8 (1.829 vs 1.829).

\section{Result Summary}

The opening experiments confirmed: the gated delta at hl$=$512 is the strongest substrate primitive. It reached 1.780 bpb at scale 12, beating the leader (1.730 at scale 8 with routing collapse, more parameters, and 175k TFLOPs) in compute efficiency. But 1.780 was not 1.4. The delta is a selective writer, not an order encoder. The gate fires on difficulty ($r = -0.65$), not position.


\chapter{Fleet Infrastructure}

\section{Four Critical Kubernetes Bugs}

The fleet ran on 3 GPU hosts (slop-01, slop-02 with A4000 16GB each; slop-home with Quadro RTX 4000 8GB) under Kubernetes. Four bugs were discovered and fixed during the first 12 hours:

\subsection{Bug 1: Name Reuse Race}

K8s job names were deterministic: \texttt{chronohorn-<run-name>}. The sequence \texttt{stop\_run()} (which does async delete) followed by immediate re-dispatch created a race: the new job was applied before the old one finished deleting. K8s saw a conflict and killed the new job. Running experiments were lost.

\textbf{Fix:} Synchronous delete before apply. Wait for the old pod to confirm termination before submitting the new job manifest.

\subsection{Bug 2: Stale Detection on Missing Phase}

The fleet reconciliation treated ``missing'' k8s phase (returned when the API query fails transiently) as ``stale.'' Transient k8s query failures marked running jobs as stale, triggering auto-delete. Running experiments were killed by infrastructure hiccups.

\textbf{Fix:} Only ``failed'' phase triggers stale detection. ``Missing'' phase is treated as ``unknown'' and the job is left alone.

\subsection{Bug 3: Zombie Reconciliation}

When a user called \texttt{stop\_run()} manually, the DB entry was set to ``stopped.'' But reconciliation ran on a 60-second loop and overwrote the status back to ``running'' if the k8s pod was still in its termination grace period.

\textbf{Fix:} 2-minute grace window after any manual status change. Reconciliation skips jobs whose last event in the events table is $< 120$ seconds old.

\subsection{Bug 4: Auto-delete After Pull}

Completed k8s jobs remain on the cluster indefinitely. The GPU allocation check (\texttt{alloc\_gpus()}) counts all pods on a node, including completed ones. A node with 3 completed jobs shows 3 GPUs ``in use'' even though they are idle. New dispatch was blocked.

\textbf{Fix:} Delete the k8s job immediately after successful result pull. The results live in the DB; the pod is no longer needed.

\section{TTL Durable Storage}

K8s pods write to ephemeral storage that vanishes on pod termination. The fleet drain path pulls results from the pod's export directory via \texttt{kubectl cp}. If the pod terminates before drain completes, the results are lost.

\textbf{Solution:} Every training job copies ALL artifacts (result JSON, checkpoint, training state, probes) to \texttt{/data/chronohorn/checkpoints/<run-name>/} before the pod exits. This is a host-path volume mount that survives pod termination. The drain path pulls from this location as a fallback if the pod is already gone.

Additionally: auto-ship to the Sharts external drive (\texttt{/Volumes/Sharts/heinrich/session9\_compression/}) on successful pull. 35 checkpoints were preserved. Only 1 was permanently lost (push14-tied-delta512-s8-50k), and it was superseded by a better run.

\section{Ghost Jobs}

Old manifests leave pending/running entries in the DB that block dispatch. When a manifest is replaced by a newer version with different run names, the old entries remain. They appear as ``pending'' and consume a dispatch slot, but no k8s job exists for them.

\textbf{Workaround:} Manual cleanup via \texttt{chronohorn\_reset} targeting specific run names. No automated solution yet --- the DB doesn't track which manifest created a given entry.

\section{Infrastructure Time Cost}

Of the 48-hour session, approximately 12 hours were spent on fleet infrastructure. The four k8s bugs took $\sim$4 hours to diagnose and fix. TTL durable storage took $\sim$2 hours. Ghost job cleanup was $\sim$1 hour spread across the session. The remaining $\sim$5 hours were Docker image builds (adding gcc for Triton compilation), data provisioning (byte shards for the adaptive substrate), and debugging OOM conditions.

The science-to-plumbing ratio was approximately 1:4. The experiments that mattered took 3 minutes each to configure and dispatch. The infrastructure failures took hours to diagnose.


\chapter{The Compression Sweep}

\section{Experimental Design}

All runs at scale 8, 10k steps, gated delta hl512, sp1024 tokenizer, seq\_len=512. The baseline is the delta512 configuration with 256 modes (pre-scale). Each variant changes exactly one axis.

\section{Results}

\begin{center}
\begin{tabular}{lrrrl}
\toprule
Config & bpb & tok/s & $\Delta$ bpb & Note \\
\midrule
baseline (delta512, 256m) & 1.951 & 201k & --- & Reference \\
modes128 & 1.957 & 276k & $+$0.005 & 37\% faster, noise quality loss \\
modes64 & 1.958 & 219k & $+$0.006 & 9\% faster, noise quality loss \\
osc25 (oscillation freq 25) & 1.951 & 200k & $+$0.000 & Zero signal \\
osc50 (oscillation freq 50) & 1.950 & 193k & $-$0.001 & Noise \\
local4 (local window 4) & 1.968 & 199k & $+$0.017 & Worse --- delta needs local window \\
\bottomrule
\end{tabular}
\end{center}

\section{Analysis}

\finding{Fixed-frequency oscillation is dead.} Both oscillation frequencies (25 and 50 steps) produced zero measurable signal. The oscillation multiplies each mode's decay by a sinusoidal modulator: $\alpha_t = \alpha \cdot (1 + a \sin(2\pi t / T))$. At $T = 25$ or $T = 50$ within a 512-token sequence, the modulation averages out. The substrate needs input-dependent dynamics, not prescribed frequencies.

\finding{modes128 is the efficiency sweet spot.} 128 pre-scale modes (1024 post-scale at s8) run 37\% faster than 256 modes (2048 post-scale) with only 0.005 bpb loss --- within seed variance. The scan cost is linear in mode count, so halving modes nearly halves scan time.

\killed{Fixed-frequency oscillation as a substrate mechanism. Zero signal at two frequencies.}

\finding{The local path provides essential n-gram context.} Reducing the local window from 8 to 4 tokens costs 0.017 bpb. The delta substrate cannot replace the local path's contribution. The model uses two parallel information channels: the substrate for long-range weighted averages and the local path for immediate n-gram patterns. Shrinking either one hurts.


%==============================================================
\part{Heinrich Forensics}
\label{part:forensics}
%==============================================================

\chapter{The Gate}

\section{MRI on push14-delta512-s12-50k}

Heinrich performed a full MRI capture on the delta512 model at scale 12, 50k steps, 1.780 bpb. The MRI captures substrate states, gate activations, expert routing, per-band logits, and all learned projections for 200 sequences of 512 tokens.

\subsection{Write Gate Characterization}

\begin{center}
\begin{tabular}{lrl}
\toprule
Metric & Value & Interpretation \\
\midrule
Position correlation ($r$) & 0.026 & Zero --- gate does not know position \\
Difficulty correlation ($r$) & $-$0.65 & Strong --- hard tokens close the gate \\
Effective rank & 1.2 & Collapsed to $\sim$1 dimension \\
PC1 explained variance & 97.2\% & Single dominant direction \\
Mean gate value (easy tokens) & 0.017 & Nearly closed \\
Mean gate value (hard tokens) & 0.036 & 2$\times$ open, still small \\
\bottomrule
\end{tabular}
\end{center}

\finding{The gate is a difficulty-adaptive EMA, not an order encoder.} The gate fires harder for unpredictable tokens ($r = -0.65$ with loss) and barely fires for predictable ones. It does not condition on position at all ($r = 0.026$). The gated delta makes the substrate write more when it encounters surprising content, making the running average more responsive to informative tokens.

\subsection{Dimensional Collapse}

Despite having 6144 modes (512 pre-scale $\times$ 12), the effective dimensionality of the substrate states is 7.2. The substrate has collapsed to 7 independent dimensions out of 6144 available. The gate's rank-1 structure (97.2\% in PC1) means all modes receive approximately the same gate value, differing only in their decay rate.

\subsection{Substrate vs Local Path}

\begin{center}
\begin{tabular}{lrrl}
\toprule
Path & L2 norm & \% of total & Contribution \\
\midrule
Substrate & 0.28 & 0.03\% & Context accumulation \\
Local & 920 & 99.97\% & N-gram prediction \\
\bottomrule
\end{tabular}
\end{center}

The model is 99.97\% local path by L2 magnitude. The substrate contributes 0.03\% of the signal. Yet this 0.03\% matters: removing the substrate entirely costs $>$0.1 bpb (see Chapter~\ref{ch:balance}).

\subsection{Routing Collapse (Again)}

Without bands, the delta model exhibits the same routing pathology as the leader:

\begin{center}
\begin{tabular}{lr}
\toprule
Expert & Token share \\
\midrule
Expert 0 & 100\% \\
Expert 1 & 0\% \\
Switch rate & 0\% \\
\bottomrule
\end{tabular}
\end{center}

\finding{Routing collapse is universal without bands.} Every model tested without the band architecture collapses to a single expert. The winner-take-all dynamics of softmax routing overwhelm the balance loss at the default coefficient (0.01). Bands prevent collapse by construction --- each band's experts operate on a different timescale slice of the substrate, ensuring they cannot converge.


\chapter{The Space Prefix}

\section{The Sentencepiece Encoding}

The sp1024 tokenizer prepends a special Unicode character \texttt{\_} (U+2581, ``lower one eighth block'') to tokens that appear after whitespace. Thus ``the cat'' tokenizes as \texttt{[\_the, \_cat]} where \texttt{\_the} and \texttt{the} are distinct vocabulary entries with distinct embeddings.

\section{Space-Prefix Tokens Are Harder}

\begin{center}
\begin{tabular}{lrrl}
\toprule
Token class & Mean bpb & Count & Note \\
\midrule
Bare tokens (no prefix) & $\sim$1.5 & $\sim$40\% & Continuation bytes \\
\texttt{\_} tokens (space prefix) & $\sim$4.7 & $\sim$60\% & Word boundaries \\
Ratio & 3.2$\times$ & & \\
\bottomrule
\end{tabular}
\end{center}

\finding{The cold start is every word boundary, not just position 0--4.} Each \texttt{\_} token begins a new word. The model must predict the first character of a new word from context alone --- the local window does not yet contain any characters of the new word. This is the same problem as the cold start at sequence position 0, but it occurs at every word boundary.

\section{Embedding Analysis}

\begin{center}
\begin{tabular}{lrrl}
\toprule
Pair & Embed cosine & Substrate cosine & Interpretation \\
\midrule
\texttt{\_t} vs \texttt{t} & 0.30 & 0.22 & Nearly orthogonal \\
\texttt{\_the} vs \texttt{the} & 0.28 & 0.19 & Nearly orthogonal \\
\texttt{\_a} vs \texttt{a} & 0.25 & 0.17 & Nearly orthogonal \\
\bottomrule
\end{tabular}
\end{center}

\texttt{\_t} and \texttt{t} are almost orthogonal in both embedding and substrate space. The tokenizer treats them as completely different tokens despite sharing the same characters. The model learns entirely different representations for word-initial vs word-internal occurrences.

\section{Gate Behavior at Word Boundaries}

\begin{center}
\begin{tabular}{lrl}
\toprule
Token class & Mean gate value & $\Delta$ \\
\midrule
Bare tokens & 0.017 & --- \\
\texttt{\_} tokens & 0.036 & 2.1$\times$ \\
\bottomrule
\end{tabular}
\end{center}

The gated delta writes 2$\times$ harder at word boundaries. The substrate L2 norm is 2.8$\times$ smaller for \texttt{\_} tokens despite the gate opening wider --- the gate opens, but the substrate accumulates LESS because the word-initial token embedding points in a different direction from the running average.

\finding{The tokenizer encodes local order for free.} The \texttt{\_} prefix distinguishes word-initial from word-internal positions at the vocabulary level. The substrate does not need to encode this aspect of order --- the tokenizer already did. The substrate only needs to encode inter-token order: the fact that ``A B C'' differs from ``C B A'' across word boundaries.


\chapter{Position $R^2$ Across All Models}

\section{Methodology}

For each model, we captured substrate states at all 512 positions across 200 sequences (102,400 total observations). We then fit a linear regression: $\text{position} \sim \text{substrate states}$. The corrected $R^2$ adjusts for the number of predictors to prevent overfitting artifacts.

A Heinrich bug fix corrected the $R^2$ computation: the original implementation used raw $R^2$ which inflates with mode count (at 6144 modes, even random states produce $R^2 > 0.05$). The corrected version uses adjusted $R^2$, which penalizes for the number of parameters.

\section{Position $R^2$ Results}

\begin{center}
\begin{tabular}{lrl}
\toprule
Model & Position $R^2$ & Note \\
\midrule
cb-s8-experts-50k (leader) & 0.001 & Original architecture \\
bands4-exp2-512m-s12 & 0.001 & With bands \\
push14-delta512-s12-50k & 0.001 & With gated delta \\
delta512-s8-baseline & 0.001 & Delta at scale 8 \\
complex-rotation-s8 & 0.001 & SO(2) rotation \\
lasso-s8-10k & 0.001 & GL(2) lasso rotation \\
lasso-pos-s8 & 0.002 & Lasso + position signal \\
quaternion-s8 & 0.001 & SO(3) Hamilton product \\
so5-s8 & 0.001 & SO(5) Lie algebra \\
lasso-blocks2-s8 & 0.001 & 2-block deep lasso \\
lasso-pos-blocks2-s8 & 0.003 & Lasso + position + depth \\
adaptive-s8 & 0.001 & Adaptive complex substrate \\
\bottomrule
\end{tabular}
\end{center}

\finding{Position $R^2 = 0.001$ across ALL 12 models tested. Zero. Definitive.} No rotation algebra, no depth, no position signal moves the needle. The substrate does not encode position. This is the single most important finding of Session 9.

\section{Content $R^2$ Results}

Content $R^2$ measures how well the substrate states predict the token at each position (a classification probe for the 1024-token vocabulary). This is the substrate's ``content capacity'' --- how much information about individual tokens survives the exponential averaging.

\begin{center}
\begin{tabular}{lrrl}
\toprule
Model & Content $R^2$ & MLP content $R^2$ & $\Delta$ vs baseline \\
\midrule
baseline (no rotation) & 0.101 & 0.209 & --- \\
SO(2) complex & 0.105 & 0.215 & $+$0.004 (noise) \\
\textbf{Lasso GL(2)} & \textbf{0.209} & \textbf{0.321} & $+$\textbf{0.108} \\
Quaternion SO(3) & 0.245 & 0.344 & $+$0.144 \\
SO(5) Lie algebra & 0.228 & 0.346 & $+$0.127 \\
Lasso + pos + blocks2 & 0.217 & 0.332 & $+$0.116 \\
\bottomrule
\end{tabular}
\end{center}

\finding{Every mechanism that improved bpb did so by increasing content $R^2$, not position $R^2$.} The lasso doubled content $R^2$ from 0.101 to 0.209. Quaternion and SO(5) pushed it further. But position $R^2$ did not move. The improvement comes from richer content coupling --- more dimensions of WHAT tokens meant --- not from encoding WHEN tokens appeared.

\section{The MLP Content $R^2$ Ceiling}

The MLP content $R^2$ (a 2-layer MLP probe instead of linear) saturates at $\sim$0.35 across all noncommutative algebras (quaternion: 0.344, SO(5): 0.346, lasso: 0.321). This ceiling means 65\% of prediction comes from the local path and embedding, not the substrate. Breaking this ceiling requires a fundamentally different readout architecture, not a better rotation.


%==============================================================
\part{The Rotation}
\label{part:rotation}
%==============================================================

\chapter{Complex Rotation --- SO(2)}

\section{The Idea}

The simplest possible rotation: multiply each mode pair by a unit complex number before decaying.
\[
z_t = \alpha \cdot e^{i\omega(x_t)} \cdot z_{t-1} + g(x_t) \cdot v(x_t)
\]
where $\omega(x_t) = W_\omega \cdot x_t$ is an input-dependent phase angle. The $e^{i\omega}$ factor rotates the complex state by an angle that depends on the current token. Geometrically, each token ``twists'' the substrate state by a learned angle.

\section{Why SO(2) Should Work}

The rotation breaks the EMA's shift invariance. Without rotation, the recurrence $h_t = \alpha h_{t-1} + v_t$ produces a weighted sum that is invariant to input reordering (up to the exponential decay). With rotation, the accumulated phase $\sum_{k=0}^{t} \omega(x_k)$ creates an input-dependent state that remembers the total ``twist'' --- a summary of all tokens seen so far.

\section{Why SO(2) Can't Work}

Complex multiplication commutes: $e^{i\omega_A} \cdot e^{i\omega_B} = e^{i(\omega_A + \omega_B)} = e^{i\omega_B} \cdot e^{i\omega_A}$. The accumulated phase is a SUM of per-token angles. Sums don't remember order. The sequence ``A B C'' and ``C B A'' produce the same accumulated phase. SO(2) rotation is commutative.

\section{Implementation}

The initial implementation was a sequential Python for-loop over sequence positions. Throughput: 8--12k tok/s. Unusable at scale.

The fix: complex parallel prefix scan. The recurrence $z_t = a_t z_{t-1} + b_t$ (where $a_t = \alpha \cdot e^{i\omega_t}$ and $b_t = g_t \cdot v_t$) is an associative scan with the operator:
\[
(a_1, b_1) \circ (a_2, b_2) = (a_1 \cdot a_2,\ a_2 \cdot b_1 + b_2)
\]
This parallelizes with Blelloch/Hillis-Steele scan in $O(\log n)$ depth. After the fix: 200k tok/s. A 20$\times$ speedup from algorithmic parallelism.

\section{Results}

\begin{center}
\begin{tabular}{lrrrl}
\toprule
Config & bpb & tok/s & Params & Note \\
\midrule
Baseline (no rotation) & 1.951 & 201k & --- & --- \\
Complex SO(2) & 1.943 & 200k & 4k & $-$0.008 (noise) \\
\bottomrule
\end{tabular}
\end{center}

\subsection{Heinrich MRI}

The complex rotation changed the substrate's internal dynamics substantially even though the bpb improvement was negligible:

\begin{center}
\begin{tabular}{lrrl}
\toprule
Metric & Baseline & SO(2) & $\Delta$ \\
\midrule
Gate opening (mean) & 2.5\% & 19.7\% & 8$\times$ \\
Content variance concentration & low & high & Gate focuses writes \\
Loss correlation with gate & 0.225 & 0.351 & Gate is more predictive \\
Position $R^2$ & 0.001 & 0.001 & Unchanged \\
Content $R^2$ & 0.101 & 0.105 & Unchanged \\
\bottomrule
\end{tabular}
\end{center}

The gate opened 8$\times$ wider, content variance concentrated, and the gate became more correlated with loss. But position $R^2$ did not move. The accumulated phase is a sum --- sums don't remember order.

\killed{SO(2) complex rotation. Commutative algebra cannot encode order. $-$0.008 bpb is noise. The gate dynamics changed but the downstream bpb didn't follow.}


\chapter{The Lasso --- GL(2)}

\section{The Key Insight}

Scalars commute. Complex numbers commute. $2 \times 2$ matrices do NOT commute:
\[
\begin{pmatrix} a & b \\ c & d \end{pmatrix} \begin{pmatrix} e & f \\ g & h \end{pmatrix}
\neq
\begin{pmatrix} e & f \\ g & h \end{pmatrix} \begin{pmatrix} a & b \\ c & d \end{pmatrix}
\]

The lasso replaces the scalar (or complex) decay with a $2 \times 2$ matrix transition per mode pair:
\[
h_t = M(x_t) \cdot h_{t-1} + g(x_t) \cdot v(x_t)
\]
where $M(x_t) \in \text{GL}(2)$ is an input-dependent $2 \times 2$ matrix (the general linear group --- invertible matrices, no norm constraint). The noncommutativity of matrix multiplication means ``A B'' and ``B A'' produce different states. The substrate can distinguish input order.

The name ``lasso'' comes from the topological interpretation: a $2 \times 2$ matrix acts on a 2D plane, and the composition of input-dependent rotations traces a path (a lasso) in GL(2) that depends on the order of inputs. Different orderings trace different lassos.

\section{Parallel Prefix Scan}

The recurrence $h_t = M_t \cdot h_{t-1} + b_t$ is an associative scan with:
\[
(M_1, b_1) \circ (M_2, b_2) = (M_1 \cdot M_2,\ M_2 \cdot b_1 + b_2)
\]

The implementation uses Hillis-Steele scan with \texttt{torch.roll} + \texttt{torch.where} instead of explicit indexing. This is compatible with \texttt{torch.compile} (dynamo):

\begin{verbatim}
for d in range(int(math.ceil(math.log2(T)))):
    shift = 1 << d
    M_shifted = torch.roll(M, shift, dims=-3)
    b_shifted = torch.roll(b, shift, dims=-2)
    mask = (positions >= shift)
    M = torch.where(mask, M_shifted @ M, M)
    b = torch.where(mask, M_at_t @ b_shifted + b, b)
\end{verbatim}

The scan runs in $O(n \log n)$ time and $O(n)$ space, with $\log_2(512) = 9$ passes for a 512-token sequence. Each pass is a batched matrix multiply over the entire sequence, fully utilizing GPU parallelism.

\section{Results}

\begin{center}
\begin{tabular}{lrrrrl}
\toprule
Config & bpb & tok/s & Params & $\Delta$ bpb & Note \\
\midrule
Baseline & 1.951 & 201k & --- & --- & --- \\
SO(2) complex & 1.943 & 200k & 4k & $-$0.008 & Commutative \\
\textbf{Lasso GL(2)} & \textbf{1.842} & \textbf{316k} & \textbf{16k} & $-$\textbf{0.109} & \textbf{Noncommutative} \\
\bottomrule
\end{tabular}
\end{center}

\finding{The lasso produces $-$0.109 bpb, the largest legal improvement at matched speed.} The improvement is 14$\times$ larger than SO(2) ($-$0.109 vs $-$0.008). The lasso adds only 16,448 parameters --- 0.2\% of the model. Seed variance is $\pm$0.0001, making the signal 1,000$\times$ the noise.

\section{Heinrich Forensics}

\begin{center}
\begin{tabular}{lrrl}
\toprule
Metric & Baseline & Lasso & $\Delta$ \\
\midrule
Content $R^2$ (linear) & 0.101 & 0.209 & $+$0.108 (doubled) \\
Content $R^2$ (MLP) & 0.209 & 0.321 & $+$0.112 \\
Effective dimensionality & 81.6 & 122.0 & $+$40.4 \\
Substrate L2 norm & 0.13 & 0.29 & 2.2$\times$ \\
Position $R^2$ & 0.001 & 0.001 & Zero \\
\bottomrule
\end{tabular}
\end{center}

\finding{The lasso doubled content $R^2$ from 0.101 to 0.209.} The substrate now contains twice as much information about individual tokens at each position. The effective dimensionality exploded from 81.6 to 122.0 --- the $2 \times 2$ coupling created 40 new independent dimensions for the readout to exploit.

\finding{Position $R^2$ is still 0.001.} The lasso's improvement comes entirely from richer content coupling, not order encoding. The substrate is a better accumulator of WHAT tokens appeared, not WHEN they appeared.

\section{Why the Lasso Helps Without Encoding Order}

The $2 \times 2$ matrix transition couples mode pairs. In the scalar EMA, each mode evolves independently: $h_t^{(k)} = \alpha_k h_{t-1}^{(k)} + v_t^{(k)}$. In the lasso, pairs of modes interact: the matrix $M_t$ mixes mode $k$ with mode $k+1$ at every step. This cross-mode coupling creates new representational dimensions that the readout can use.

Consider two modes with half-lives 5 and 50. In the scalar EMA, they independently track fast and slow averages. In the lasso, the fast mode's state bleeds into the slow mode and vice versa. The result is not just ``fast average'' and ``slow average'' but a richer manifold that captures how fast-scale patterns relate to slow-scale patterns. The readout can now distinguish contexts that differ in their fast-slow correlation, not just their fast and slow averages independently.

The noncommutativity is necessary for this coupling to be input-dependent. If the matrices commuted, the coupling would be the same regardless of input order, and the manifold would collapse back to independent averages. The lasso's noncommutativity creates input-dependent coupling patterns --- richer content, not position.


\chapter{The Cayley-Dickson Hierarchy}

\section{Overview}

The Cayley-Dickson construction generates algebras of increasing dimension:
\[
\mathbb{R} \xrightarrow{\text{double}} \mathbb{C} \xrightarrow{\text{double}} \mathbb{H} \xrightarrow{\text{double}} \mathbb{O} \xrightarrow{\text{double}} \mathbb{S} \xrightarrow{\text{double}} \cdots
\]

At each doubling, the algebra gains dimension but loses a structural property: commutativity at $\mathbb{H}$ (quaternions), associativity at $\mathbb{O}$ (octonions), alternativity at $\mathbb{S}$ (sedenions). We tested representatives at each level plus unconstrained matrix groups of larger dimension.

\section{Full Results Table}

\begin{center}
\begin{tabular}{lrrrrrrl}
\toprule
Algebra & Dim & Comm. & Assoc. & bpb & tok/s & Params & Note \\
\midrule
Real (scalar EMA) & 1 & yes & yes & 1.951 & 201k & 0 & Baseline \\
Complex SO(2) & 2 & yes & yes & 1.943 & 200k & 4k & Commutative \\
\textbf{Lasso GL(2)} & \textbf{4} & \textbf{no} & \textbf{yes} & \textbf{1.842} & \textbf{316k} & \textbf{16k} & \textbf{Best speed/bpb} \\
Quaternion SO(3) & 4 & no & yes & 1.844 & 293k & 8k & Matches lasso \\
Quintic GL(5) & 25 & no & yes & NaN & --- & 51k & Diverged \\
SO(5) Lie exp & 10 & no & yes & 1.845 & 118k & 21k & Stable, 2.7$\times$ slower \\
\bottomrule
\end{tabular}
\end{center}

\section{Quaternion SO(3)}

The quaternion rotation uses the Hamilton product to compose 4D rotations:
\[
q = a + bi + cj + dk, \qquad q_1 \cdot q_2 = (a_1 a_2 - b_1 b_2 - c_1 c_2 - d_1 d_2) + \cdots
\]
Unit quaternions ($|q| = 1$) parameterize SO(3) --- the group of 3D rotations. Non-unit quaternions include scaling, giving GL(1, $\mathbb{H}$).

\textbf{Result:} 1.844 bpb, matching the lasso (1.842) within seed variance. But 7\% slower at 293k tok/s vs 316k tok/s. The quaternion product requires more operations per step than the $2 \times 2$ matrix multiply.

We tested with and without norm constraint ($|q| = 1$ vs unconstrained): the norm constraint does not help. Unconstrained quaternions (1.844) match unit quaternions (1.846). The spectral radius of the quaternion naturally stays bounded because the gate and decay provide sufficient regularization.

\section{Quintic GL(5)}

An unconstrained $5 \times 5$ matrix transition per mode group. 25 parameters per mode group, 51k total.

\textbf{Result:} NaN. The model diverged within 500 training steps.

\textbf{Why:} The spectral radius of a product of unconstrained $5 \times 5$ matrices grows exponentially with sequence length. In a 512-step prefix scan, the product of 512 matrices has spectral radius $\gg 1$. The state magnitude explodes, producing NaN logits.

The $2 \times 2$ lasso does not diverge because the gate ($g \approx 0.02$) and decay ($\alpha = 0.9987$) jointly constrain the spectral radius. At $5 \times 5$, the degrees of freedom are sufficient for the product to find unstable directions that escape the gate's control.

\section{SO(5) via Lie Algebra}

To fix the divergence, we constrained the $5 \times 5$ matrix to SO(5) --- the special orthogonal group. Every element of SO(5) has spectral radius exactly 1 (eigenvalues on the unit circle), so the state cannot diverge.

The parameterization: project from the input to the Lie algebra $\mathfrak{so}(5)$ (the space of $5 \times 5$ antisymmetric matrices, dimension 10), then exponentiate:
\[
M(x_t) = \exp(A(x_t)), \qquad A = -A^T \in \mathfrak{so}(5)
\]
The matrix exponential guarantees $M \in \text{SO}(5)$ regardless of the input $A$. Stability by construction.

\textbf{Result:} 1.845 bpb, matching the lasso (1.842) within seed variance. But 2.7$\times$ slower at 118k tok/s, due to the \texttt{torch.linalg.matrix\_exp} call at every step. The matrix exponential computes a Pad\'{e} approximation, which is expensive for batched $5 \times 5$ matrices.

\textbf{Requires \texttt{--state-dim 40 --num-heads 8}} (head\_dim $= 5$) to align the state dimensions with SO(5)'s 5-dimensional representation.

\section{The Key Finding}

\finding{Non-solvability doesn't matter. GL(2), SO(3), and SO(5) all converge to $\sim$1.844 bpb.} SO(5) is the first non-solvable compact group in the hierarchy \cite{abel} --- its Lie algebra $\mathfrak{so}(5)$ has no composition series of abelian quotients. GL(2) and SO(3) are solvable. The distinction, which is central to Galois theory and the unsolvability of the quintic \cite{cayley}, makes no difference to language modeling. The improvement comes from noncommutativity itself, not the specific algebraic structure.

\finding{Depth doesn't stack.} Testing deeper compositions:

\begin{center}
\begin{tabular}{lrrl}
\toprule
Config & bpb & $\Delta$ vs lasso & Note \\
\midrule
Lasso (1 block) & 1.842 & --- & Baseline \\
Lasso blocks2 & 1.844 & $+$0.002 & Noise \\
Lasso blocks4 & 1.850 & $+$0.008 & Worse \\
\bottomrule
\end{tabular}
\end{center}

Content coupling saturates in one pass. The readout can extract all the benefit from a single layer of noncommutative coupling. Adding more layers does not provide additional useful dimensions.


\chapter{Position Signal}

\section{Breaking Shift Invariance}

The position signal feeds $\log(1 + t)$ into the gate and lasso projections:
\[
\text{input}_t = [x_t;\ \log(1 + t)]
\]
This explicitly breaks shift invariance --- the model's behavior at position 100 differs from position 200 even for identical token sequences. The logarithmic scaling ensures the signal grows slowly, avoiding domination at late positions.

\section{Results}

\begin{center}
\begin{tabular}{lrrrl}
\toprule
Config & bpb & tok/s & $\Delta$ & Note \\
\midrule
Lasso (no position) & 1.842 & 316k & --- & --- \\
Lasso + position & 1.841 & 313k & $-$0.001 & Noise \\
Lasso + position + blocks2 & 1.836 & 308k & $-$0.006 & Best result \\
\bottomrule
\end{tabular}
\end{center}

\section{Gate Position Correlation}

\begin{center}
\begin{tabular}{lrrl}
\toprule
Model & Gate--position $r$ & bpb & Note \\
\midrule
Lasso (no position) & 0.026 & 1.842 & Gate ignores position \\
Lasso + position & $-$0.71 & 1.841 & Gate strongly conditions on position \\
\bottomrule
\end{tabular}
\end{center}

\finding{The model HAS access to position but doesn't USE it at 10k steps.} The gate strongly conditions on position ($r = -0.71$) when the signal is provided, but bpb doesn't improve (1.841 vs 1.842). The gate adjusts its write strength based on position --- writing harder at early positions and softer at late positions --- but this adaptive write schedule doesn't improve prediction.

\section{The Position + Depth Interaction}

Individually, position signal ($-$0.001) and depth ($+$0.002) are both noise. Together, they stack:

\begin{center}
\begin{tabular}{lrrl}
\toprule
Config & bpb & $\Delta$ vs lasso & Interpretation \\
\midrule
Lasso & 1.842 & --- & --- \\
Lasso + position & 1.841 & $-$0.001 & Noise alone \\
Lasso + blocks2 & 1.844 & $+$0.002 & Noise alone \\
Lasso + position + blocks2 & 1.836 & $-$0.006 & Signal together \\
\bottomrule
\end{tabular}
\end{center}

The combination costs 1\% speed (308k vs 316k tok/s) for $-$0.006 bpb. The position signal gives the second block a reference frame --- ``this is an early-position refinement'' vs ``this is a late-position refinement.'' Without the position signal, the second block has nothing to differentiate its behavior from the first.

\alive{Position + depth as a combined mechanism. 1.836 bpb, best at scale 8.}


%==============================================================
\part{The Speed Story}
\label{part:speed}
%==============================================================

\chapter{From 8k to 328k tok/s}

\section{The Speed Timeline}

\begin{center}
\begin{tabular}{lrrl}
\toprule
Configuration & tok/s & Improvement & Change \\
\midrule
Sequential Python loop (SO(2)) & 8--12k & baseline & Initial implementation \\
Parallel prefix scan (roll+where) & 200k & 20$\times$ & Algorithmic parallelism \\
+ batch=32 & 235k & 1.2$\times$ & Better GPU utilization \\
+ torch.compile (default mode) & 295k & 1.5$\times$ & Kernel fusion \\
+ batch=64 & 316k & 1.6$\times$ & Max VRAM utilization \\
\bottomrule
\end{tabular}
\end{center}

The total improvement from the initial sequential implementation to the final optimized version is \textbf{30--40$\times$}. The two dominant factors are:

\begin{enumerate}[nosep]
\item \textbf{Algorithmic parallelism (20$\times$):} The sequential loop processes one position at a time with $O(n)$ depth. The Hillis-Steele scan processes all positions simultaneously in $O(\log n)$ depth. For $n = 512$, this is $512 / 9 \approx 57\times$ theoretical parallelism. The practical gain is 20$\times$ after memory bandwidth and kernel launch overhead.

\item \textbf{Batching + compilation (1.6$\times$):} Larger batch sizes provide more tokens per matmul, better utilizing the GPU's tensor cores. \texttt{torch.compile} fuses elementwise operations (roll, where, multiply) into a single kernel launch, reducing overhead.
\end{enumerate}

\section{torch.compile Modes}

\begin{center}
\begin{tabular}{lrl}
\toprule
Mode & Result & Note \\
\midrule
default (kernel fusion) & Works & 1.5$\times$ speedup \\
max-autotune & Crashes & CUDA graph tensor aliasing in scan \\
reduce-overhead & Crashes & Same CUDA graph issue \\
\bottomrule
\end{tabular}
\end{center}

Both \texttt{max-autotune} and \texttt{reduce-overhead} attempt to capture the scan as a CUDA graph. The Hillis-Steele scan modifies tensors in-place across iterations (\texttt{torch.roll} + \texttt{torch.where} with same output tensor), which violates CUDA graph's static allocation assumption. The tensors alias across graph replay, producing incorrect results or segfaults.

\textbf{Fix:} Use only default mode (kernel fusion without CUDA graphs). Additionally, add \texttt{.clone()} on the forward output to break potential aliasing for downstream CUDA graph consumers.

\section{Docker Requirements}

The Triton compiler (used by \texttt{torch.compile}) requires \texttt{gcc} at compile time. The base PyTorch Docker image does not include \texttt{gcc}. Training jobs crashed silently with ``Triton compilation failed'' until \texttt{gcc} was added to the Docker image.

\section{Auto-Scaling Defaults}

Two auto-scaling behaviors were established and are now defaults:

\begin{enumerate}[nosep]
\item \textbf{Batch size:} Auto-scales to fill 60\% of GPU VRAM. A probe run of 10 forward passes at batch=1 measures per-sample memory. The batch size is set to $\lfloor 0.6 \times \text{VRAM} / \text{per\_sample} \rfloor$. The 60\% threshold leaves headroom for activation checkpointing, optimizer states, and Triton workspace.

\item \textbf{torch.compile:} Auto-enables for training runs $\geq$ 5000 steps. Compilation adds $\sim$60 seconds of overhead on the first forward pass. For short runs ($< 5000$ steps), the overhead exceeds the total speedup. For long runs, the 1.5$\times$ sustained speedup dominates.
\end{enumerate}


\chapter{GPU Physics}

\section{A4000 Specifications}

\begin{center}
\begin{tabular}{lr}
\toprule
Metric & Value \\
\midrule
FP32 peak & 19.2 TFLOPS \\
TF32 peak (tensor cores) & 77 TFLOPS \\
Memory bandwidth & 448 GB/s \\
VRAM & 16 GB GDDR6 \\
Tensor cores & 192 (3rd gen) \\
\bottomrule
\end{tabular}
\end{center}

\section{Utilization Analysis}

At the lasso's operating point (316k tok/s, batch=64, scale 8):

\begin{center}
\begin{tabular}{lrrl}
\toprule
Component & TFLOPS & \% of TF32 peak & Bottleneck \\
\midrule
Readout (expert MLP) & 8.2 & 10.6\% & Compute \\
Substrate scan (roll+where) & 2.1 & 2.7\% & Memory BW \\
Local path (8-token conv) & 1.9 & 2.5\% & Compute \\
Embedding + projection & 1.5 & 1.9\% & Memory BW \\
\midrule
Total & 13.65 & 17.7\% & --- \\
\bottomrule
\end{tabular}
\end{center}

The model achieves 17.7\% of TF32 peak. The remaining 82\% is lost to:

\begin{enumerate}[nosep]
\item \textbf{Small matmuls:} At scale 8 with 2 experts, the readout matmuls are $[64 \times 512, 512 \times 256]$ --- too small to saturate tensor cores. The A4000 needs matmuls of dimension $\geq$2048 for 80\%+ utilization.

\item \textbf{Memory-bound scan:} The substrate scan (\texttt{torch.roll} + \texttt{torch.where}) is pure element-wise FP32. It does not use tensor cores at all. At 2048 modes $\times$ 512 positions $\times$ batch 64, the scan moves $\sim$500 MB per step, bounded by the 448 GB/s memory bandwidth.

\item \textbf{Kernel launch overhead:} The compiled scan still requires 9 kernel launches (one per Hillis-Steele pass). Each launch has $\sim$5\,$\mu$s overhead. At 316k tok/s, each step takes $\sim$100\,$\mu$s, so kernel launches consume $\sim$45\% of step time.
\end{enumerate}

\section{The Scale Barrier}

\finding{At scale 8, the model is memory-bound, not compute-bound.} The readout matmuls are too small for tensor core saturation. Increasing scale (larger matmuls) would improve utilization but increases parameter count. The fundamental tension: O($n$) scan models produce small intermediate representations (mode count $\times$ batch) that don't map well to modern GPU architectures optimized for large matmuls.

The fix is either (a) much larger batch sizes (batch=256+ to make matmuls wider), (b) fused scan kernels (mamba-ssm style \cite{gu} that eliminate per-step kernel launches), or (c) hardware that is better at small matmuls (e.g., Cerebras or Graphcore).


%==============================================================
\part{The Ablation Matrix}
\label{part:ablation}
%==============================================================

\chapter{Substrate/Local Balance}
\label{ch:balance}

\section{Removing the Local Path}

\begin{center}
\begin{tabular}{lrrrrl}
\toprule
Config & bpb & tok/s & Steps & $\Delta$ vs with-local & Note \\
\midrule
Lasso s8 (with local) & 1.842 & 316k & 10k & --- & Reference \\
Lasso s8 nolocal & 1.948 & 305k & 10k & $+$0.106 & 11\% faster \\
Lasso s12 nolocal 50k & 1.873 & --- & 50k & --- & At scale \\
Lasso s12 (with local) 50k & 1.739 & 185k & 50k & --- & Frontier \\
\bottomrule
\end{tabular}
\end{center}

\section{Analysis}

\finding{The substrate alone (no local path) at scale 8 matches the baseline WITH local path.} The nolocal s8 model (1.948) is within 0.003 bpb of the baseline-with-local (1.951). The lasso substrate has fully absorbed what the local path provides at scale 8 with 10k steps. But this parity breaks at scale:

At scale 12 with 50k steps, the local path provides \textbf{0.134 bpb} of n-gram prediction the substrate cannot replicate (1.873 nolocal vs 1.739 with-local). The local path's 8-token convolution window captures immediate byte transitions that the exponential averaging destroys, even with the lasso's noncommutative coupling.

\section{Adjusting the Local Scale}

\begin{center}
\begin{tabular}{lrrl}
\toprule
Local scale & bpb & $\Delta$ vs default (0.25) & Note \\
\midrule
0.0 (no local) & 1.948 & $+$0.106 & Substrate only \\
0.1 & 1.851 & $+$0.009 & Worse \\
0.25 (default) & 1.842 & --- & Reference \\
0.50 & 1.832 & $-$0.010 & Slight gain \\
\bottomrule
\end{tabular}
\end{center}

Local scale 0.5 provides a small $-$0.010 bpb improvement. The default (0.25) is near-optimal. More local helps slightly but at diminishing returns.

\finding{The local path provides 0.134 bpb of n-gram prediction at scale 12 that the substrate cannot replicate.} The two paths are complementary, not redundant. The substrate provides context (what has appeared over the last $\sim$500 tokens), the local path provides transition (what byte follows what byte in the immediate 8-token window).


\chapter{Readout and State Capacity}

\section{Expert Count}

\begin{center}
\begin{tabular}{lrrl}
\toprule
Experts & bpb & $\Delta$ vs 2 experts & Note \\
\midrule
2 (default) & 1.842 & --- & Reference \\
4 & 1.843 & $+$0.001 & Noise \\
\bottomrule
\end{tabular}
\end{center}

Four experts provide zero benefit over two. This is consistent with the routing collapse findings: without bands, additional experts collapse to the same computation. With bands, more experts per band might help, but we did not test this combination in Session 9.

\section{Band Count}

\begin{center}
\begin{tabular}{lrrl}
\toprule
Bands & bpb & $\Delta$ vs no bands & Note \\
\midrule
0 (default, no bands) & 1.842 & --- & Reference \\
4 & 1.852 & $+$0.010 & Worse \\
\bottomrule
\end{tabular}
\end{center}

\finding{Bands don't help with the lasso substrate.} In Session 6, bands improved bpb by 0.018 because they forced timescale differentiation in the readout. The lasso already creates cross-timescale coupling via its $2 \times 2$ matrix mixing. The bands' forced separation actually HURTS by preventing the readout from combining fast and slow information freely.

\section{State Dimension}

\begin{center}
\begin{tabular}{lrrrl}
\toprule
state\_dim & 2$\times$2 pairs & bpb & $\Delta$ & Note \\
\midrule
16 & 8 & 1.854 & $+$0.019 & Too few pairs \\
32 (default) & 16 & 1.842 & --- & Reference \\
64 & 32 & 1.830 & $-$0.005 & Slight gain \\
\bottomrule
\end{tabular}
\end{center}

State\_dim 64 (32 mode pairs) provides a marginal $-$0.005 bpb gain. The lasso's content coupling is nearly saturated at 16 pairs (state\_dim 32). Doubling the pairs provides more independent cross-timescale interactions, but the readout cannot extract much additional information from them.

\finding{32 is the saturation point for state\_dim with the lasso.} Below 32 (state\_dim 16), there are too few $2 \times 2$ pairs and the content coupling is insufficient. Above 32, gains are marginal.


\chapter{The Frontier Push}

\section{lasso-pos-blocks2 at Scale 12, 50k Steps}

The best configuration from the ablation matrix --- lasso + position signal + 2 blocks --- was trained to scale 12 for 50k steps on slop-02.

\begin{center}
\begin{tabular}{lrrrrr}
\toprule
Model & bpb & tok/s & Params & TFLOPs & Note \\
\midrule
Leader (cb-s8-experts-50k) & 1.730 & 65k & 17.3M & 175k & Routing collapsed \\
\textbf{Lasso-pos-blocks2-s12-50k} & \textbf{1.739} & \textbf{185k} & \textbf{20.1M} & \textbf{71k} & \textbf{Clean routing} \\
\bottomrule
\end{tabular}
\end{center}

\section{Comparison with the Leader}

The lasso model:
\begin{itemize}[nosep]
\item \textbf{Matches the leader:} 1.739 vs 1.730 bpb. Within 0.009 bpb, likely reachable with continued training.
\item \textbf{3$\times$ faster:} 185k vs 65k tok/s. The scan path is fundamentally faster than the kernel path.
\item \textbf{2.5$\times$ more compute-efficient:} 71k vs 175k TFLOPs to reach comparable quality.
\item \textbf{Clean routing:} No routing collapse. All experts active. Healthy gradient flow.
\end{itemize}

\section{Learning Curve}

The slope at step 50k was 0.002 bpb/10k-steps --- still learning. Extrapolating:

\begin{center}
\begin{tabular}{lrl}
\toprule
Steps & Projected bpb & Note \\
\midrule
50k (actual) & 1.739 & Measured \\
75k (extrapolated) & $\sim$1.73 & Matches leader \\
100k (extrapolated) & $\sim$1.70 & New frontier \\
\bottomrule
\end{tabular}
\end{center}

At 100k steps, the lasso model would likely reach $\sim$1.70 bpb --- a new legal frontier. The slope suggests the model has not saturated. The saturation forecast (3-parameter logistic fit) projects an asymptote of $\sim$1.62 bpb for this configuration.


%==============================================================
\part{The Adaptive Substrate}
\label{part:adaptive}
%==============================================================

\chapter{Five Projections, One Recurrence}

\section{Motivation}

The causal bank architecture has accumulated many prescribed structures over 9 sessions: fixed half-lives, bands, experts, local window, sticky registers, gated delta, lasso rotation. Each was designed by a human-AI collaboration to solve a specific problem. The adaptive substrate asks: what if we remove all prescribed structure and let the model discover everything?

\section{The Architecture}

The minimal recurrence with complex rotation and all-learned projections:
\[
z_k(t) = \underbrace{\sigma(W_{\text{decay}} \cdot x_t)}_{\text{decay}} \cdot \underbrace{e^{i \cdot W_\omega \cdot x_t}}_{\text{rotation}} \cdot z_k(t-1) + \underbrace{\sigma(W_{\text{write}} \cdot x_t)}_{\text{gate}} \cdot \underbrace{(W_{\text{proj\_re}} \cdot x_t + i \cdot W_{\text{proj\_im}} \cdot x_t)}_{\text{value}}
\]

Five projections from embed\_dim to n\_modes:
\begin{enumerate}[nosep]
\item $W_{\text{decay}}$: learned decay per mode per token (replaces fixed half-lives)
\item $W_\omega$: learned rotation angle per mode per token (replaces fixed or lasso rotation)
\item $W_{\text{write}}$: learned write gate per mode per token (replaces gated delta)
\item $W_{\text{proj\_re}}$: real part of value projection
\item $W_{\text{proj\_im}}$: imaginary part of value projection
\end{enumerate}

No local path. No bands. No heads. No blocks. No half-lives. The model discovers everything from data.

\section{Results}

\begin{center}
\begin{tabular}{lrrrl}
\toprule
Model & bpb & tok/s & Params & Note \\
\midrule
Lasso s8 10k & 1.842 & 316k & 16k (rotation) & Prescribed structure \\
\textbf{Adaptive s8 10k} & \textbf{1.812} & \textbf{26k} & \textbf{---} & \textbf{All learned} \\
\bottomrule
\end{tabular}
\end{center}

\finding{The adaptive substrate (1.812) beats the lasso (1.842) by 0.030 bpb.} Fewer prescribed structures, more learned capacity. The model discovers its own decay rates, rotation angles, and write gates, and does better than any hand-designed configuration.

\section{The Speed Problem}

26k tok/s is 12$\times$ slower than the lasso (316k tok/s). The adaptive substrate uses 2048 modes with 5 independent projections, creating massive activation tensors:
\[
5 \times [\text{batch} \times \text{seq} \times \text{modes}] = 5 \times [64 \times 512 \times 2048] = 335M \text{ floats} = 1.3 \text{ GB}
\]

The projections are dense matrix multiplies from embed\_dim (256) to n\_modes (2048), executed 5 times per token. The scan itself is O($n$) but the projection overhead dominates.

\section{Speed Optimization Path}

Two approaches queued but not yet tested:
\begin{enumerate}[nosep]
\item \textbf{Reduce mode count:} 128 modes instead of 2048. The compression sweep showed 128 modes is near-optimal for the lasso. If the adaptive substrate can maintain its advantage at 128 modes, throughput would increase $\sim$16$\times$ to $\sim$400k tok/s.

\item \textbf{Shared projection:} Use a single projection $W_{\text{shared}}$ and derive all 5 quantities from linear combinations. This reduces 5 large matmuls to 1, with a $\sim$5$\times$ speedup.
\end{enumerate}

The adaptive substrate needs speed optimization before it is practical for fleet-scale experimentation.


\chapter{Byte-Level}

\section{The Idea}

Remove the tokenizer entirely. Vocabulary = 256 (raw UTF-8 bytes). The model IS the tokenizer --- it discovers word boundaries, character patterns, and language structure from raw bytes.

\section{Data Pipeline}

Byte-level training requires byte-level shards. The existing sp1024 shards contain tokenized data. We built a parallel byte shard converter:

\begin{enumerate}[nosep]
\item Read sp1024 shards (tokenized sequences)
\item Decode each token sequence back to UTF-8 text using the sp1024 model
\item Re-encode as raw UTF-8 bytes
\item Pack into byte-level shards with the same format (uint8 instead of uint16)
\end{enumerate}

The converter uses all available CPU cores for parallel processing. Converting 10 GB of sp1024 shards takes $\sim$5 minutes on a 16-core machine.

Data provisioning was extended: \texttt{chronohorn data provision --variant byte} downloads and converts byte shards on fleet nodes.

\section{Results}

\begin{center}
\begin{tabular}{lrrrl}
\toprule
Model & bpb & tok/s & seq\_len & Note \\
\midrule
acb-byte-s64 & 0.888 & 88k & 64 & Sub-1.0 at byte level \\
acb-byte-s1024 & 0.866 & 59k & 1024 & More context helps \\
\bottomrule
\end{tabular}
\end{center}

\finding{Sub-1.0 bpb at byte level.} The adaptive causal bank reaches 0.888 bpb on raw UTF-8 bytes at seq\_len=64, and 0.866 at seq\_len=1024. The model learns to predict byte sequences without any tokenizer.

\section{Comparability}

Byte-level bpb is NOT directly comparable to sp1024 bpb. Individual bytes are locally predictable: within a UTF-8 multi-byte sequence, continuation bytes are heavily constrained. Within English words, letter frequencies are strongly conditioned by the preceding 2--3 letters. The byte model gets ``easy'' predictions (continuation bytes, common letter transitions) that the tokenized model never sees because the tokenizer collapses them.

A fair comparison requires measuring bits per byte of TEXT (the standard bpb metric), which accounts for the different prediction granularities. At sp1024, the model predicts $\sim$0.411 tokens per byte. At byte level, it predicts 1.0 tokens per byte. The byte model makes 2.4$\times$ more predictions per byte of text, but many of those predictions are trivially easy.

\section{Open Questions}

Heinrich raised two questions for the byte-level model:

\begin{enumerate}[nosep]
\item \textbf{Does $\omega$ develop frequency structure?} If the rotation angles form a Fourier-like basis, the adaptive substrate may be implementing a discrete Fourier transform over the byte sequence. This would be a learned equivalent of the HRR construction (Chapter~\ref{ch:hrr}).

\item \textbf{Does decay learn word boundaries at space bytes?} If the decay rate drops sharply at space bytes (0x20), the model has discovered word boundaries from data. This would be the byte-level analog of the \texttt{\_} prefix behavior observed in sp1024 (Chapter 6).
\end{enumerate}

Both questions require Heinrich MRI captures on the byte-level checkpoints, which are queued behind the speed optimization work.


\chapter{HRR and the Fourier Basis}
\label{ch:hrr}

\section{Holographic Reduced Representations}

Plate's Holographic Reduced Representations \cite{plate} use circular convolution to bind key-value pairs in a fixed-dimensional vector. The binding operation is:
\[
\text{bind}(k, v) = \mathcal{F}^{-1}[\mathcal{F}(k) \odot \mathcal{F}(v)]
\]
where $\mathcal{F}$ is the discrete Fourier transform and $\odot$ is element-wise multiplication. In the frequency domain, binding is multiplication. Retrieval is correlation (element-wise conjugate multiply).

\section{The Complex Recurrence IS Binding}

The adaptive substrate's complex recurrence has the form:
\[
z_k(t) = \alpha_k(t) \cdot e^{i\omega_k(t)} \cdot z_k(t-1) + g_k(t) \cdot v_k(t)
\]

If $\omega_k$ produces orthogonal keys (phases that don't overlap across modes), then:
\begin{itemize}[nosep]
\item The accumulated state $z_k(t)$ is a superposition of phase-rotated value vectors
\item Each token's contribution is stored at a unique phase angle per mode
\item The readout can retrieve past tokens by phase --- demodulating at the matching frequency
\end{itemize}

This IS HRR binding, implemented as a recurrent neural network.

\section{HRR Initialization}

The construction: set $\omega$ bias $= \text{linspace}(0, 2\pi, n\_\text{modes})$. Each mode gets a unique base frequency. The state vector is a DFT of the token sequence, with each mode tracking a different Fourier component.

With this initialization:
\begin{itemize}[nosep]
\item Mode $k$ oscillates at frequency $k \cdot 2\pi / n_\text{modes}$
\item The Nyquist frequency is $n_\text{modes} / 2$
\item Position information is encoded in the PHASE of each mode --- the phase at mode $k$ tells you how many tokens have passed modulo the mode's period
\end{itemize}

\section{The Frozen Omega Experiment}

The cleanest test: freeze $\omega$ at the Fourier initialization. Train only decay, write gate, value projection, and readout. The dynamics ARE the physics --- fixed Fourier basis, position from the forward pass.

If this works (position $R^2$ moves above 0.001), the construction was never needed. The existence of the Fourier basis IS the construction. The model doesn't need to learn order --- order falls out of the recurrence dynamics when the rotation frequencies are set to form a DFT.

If this fails (position $R^2$ stays at 0.001), then the HRR framework doesn't apply to language modeling. The readout cannot demodulate the Fourier components, and position information is truly inaccessible to the O($n$) substrate.

\finding{The frozen omega experiment is the decisive test for the order hypothesis. It is queued on the fleet behind speed ablation and byte curriculum.}


%==============================================================
\part{The Honest Assessment}
\label{part:honest}
%==============================================================

\chapter{What We Learned}

\section{The Eight Findings}

\begin{enumerate}
\item \textbf{The architecture wall is content capacity, not order.} Every mechanism that improved bpb did so by increasing content $R^2$, not position $R^2$. The lasso's $-$0.109 bpb doubled content $R^2$ from 0.101 to 0.209 while position $R^2$ remained at 0.001. The substrate is a content accumulator that the readout queries for ``what appeared'' not ``when it appeared.''

\item \textbf{The lasso's $-$0.109 bpb is the largest legal improvement at matched speed.} GL(2) noncommutative $2 \times 2$ matrix, 16,448 parameters, implemented as a parallel prefix scan in 9 lines of \texttt{torch.roll} and \texttt{torch.where}. The simplest noncommutative algebra that works.

\item \textbf{Position $R^2 = 0.001$ across 12+ models, 6 rotation algebras, with and without explicit position signal.} This is definitive. The substrate does not encode position regardless of what algebraic structure is imposed on the recurrence.

\item \textbf{Non-solvability doesn't matter.} GL(2), SO(3), and SO(5) all converge to $\sim$1.844 bpb. The distinction between solvable and non-solvable groups, which is central to the Cayley-Dickson hierarchy and Galois theory \cite{abel, cayley}, makes no measurable difference to language modeling.

\item \textbf{Depth doesn't stack.} Content coupling saturates in one pass. Adding blocks2 or blocks4 provides zero or negative benefit. The readout extracts all available information from a single layer of noncommutative coupling.

\item \textbf{The substrate is 0.03\% of the signal by L2.} The local path contributes 99.97\%. But the substrate's 0.03\% provides 0.109 bpb --- a disproportionate contribution relative to its magnitude. Quality is not proportional to magnitude.

\item \textbf{The adaptive substrate (1.812) beats the lasso (1.842).} Fewer prescribed structures, more learned capacity. The model discovers its own decay rates, rotation angles, and write gates, and outperforms every hand-designed configuration by 0.030 bpb.

\item \textbf{Byte-level works (0.866 bpb).} The causal bank can operate directly on raw UTF-8 bytes without a tokenizer. Sub-1.0 bpb at byte level. The model learns character patterns, word boundaries, and language structure from raw bytes.
\end{enumerate}

\section{The Content vs Order Distinction}

This is the central result of Session 9 and deserves emphasis.

The session began with the hypothesis: the 0.94 bpb locked behind temporal attention \cite{session7} is accessible if we find the right algebraic structure to encode order into the substrate. We tested six algebras across the Cayley-Dickson hierarchy, from the trivial reals through GL(5) and the Lie algebra of SO(5).

The result is unambiguous: \textbf{none of them encode order}. Position $R^2 = 0.001$ for every algebra tested. The EMA substrate, regardless of what rotation or coupling is applied, converges to a position-independent representation. The exponential decay washes out temporal information on a timescale set by the half-life, and no algebraic structure on the transition can prevent this.

But several of them improve bpb. The improvement comes from content $R^2$ --- the substrate's capacity to represent what tokens appeared, not when. The lasso's $2 \times 2$ coupling creates cross-timescale interactions that the scalar EMA cannot. The quaternion and SO(5) rotations do the same. The mechanism is dimension expansion: the noncommutative coupling creates new representational dimensions (effective dim: 81.6 $\to$ 122.0) that the readout exploits.

The distinction matters for future architecture design. If the wall were order, the solution would be a better algebra. Since the wall is content, the solution is a richer readout or a fundamentally different substrate (e.g., the adaptive substrate, which learns its own coupling structure).


\chapter{What Remains}

\section{The Open Problems}

\begin{enumerate}
\item \textbf{The 0.94 bpb from temporal attention is still untapped.} No legal mechanism has accessed the order information that temporal attention exploits. The gap between 1.739 (best legal) and 0.974 (illegal temporal attention) is 0.765 bpb. Even accounting for the causality violation inflating the temporal attention result, the true legal ceiling is likely far below 1.739.

\item \textbf{Position $R^2$ hasn't moved.} The substrate is an orderless content accumulator at every algebra tested. The frozen omega / HRR experiment is the next test: does a fixed Fourier basis provide position information through the forward-pass dynamics, without optimization?

\item \textbf{The HRR/frozen-omega experiments are queued.} The hypothesis: if $\omega$ is initialized as a DFT basis and frozen, the phase of each mode at each position encodes position information by construction. The readout can learn to demodulate. If this works, the wall was never about finding the right algebra --- it was about letting the algebra find itself.

\item \textbf{Byte-level + curriculum (64 $\to$ 4096 bytes) needs testing.} Does the model discover word boundaries? Does the decay rate develop structure at space bytes (0x20)? Byte-level removes the tokenizer bottleneck entirely, but the speed penalty (26k tok/s) must be resolved first.

\item \textbf{sp4096 + tied readout: the vocab compression lever.} The three-wall decomposition \cite{session7} attributes 62\% of the gap to vocabulary. A larger vocab (4096 vs 1024 tokens) reduces byte fallback from 0.9\% to $<$0.1\%. Blocked on data provisioning for sp4096 shards and the FLOP guard fix for tied readout.

\item \textbf{The adaptive substrate needs speed optimization.} 26k tok/s is impractical for fleet experimentation. Reducing to 128 modes and/or shared projections should bring throughput to $\sim$400k tok/s.

\item \textbf{The MLP content $R^2$ ceiling at $\sim$0.35.} 65\% of prediction comes from the local path and embedding, not the substrate. Breaking this ceiling requires a fundamentally different readout --- one that can extract more from the substrate than a 2-layer MLP can represent.
\end{enumerate}


\chapter{Infrastructure Debt}

\section{Problems Fixed in Session 9}

\begin{center}
\begin{tabular}{ll}
\toprule
Problem & Fix \\
\midrule
K8s name reuse race & Synchronous delete before apply \\
Stale detection on missing phase & Only ``failed'' triggers stale \\
Zombie reconciliation & 2-minute grace window from events \\
Auto-delete after pull & Delete k8s job after successful pull \\
TTL durable storage & Host-path volume mount survives pod exit \\
Auto-ship to Sharts drive & Copy on successful drain pull \\
FLOP guard for tied readout & Compare against embed\_dim, not vocab\_size \\
Docker gcc for Triton & Added to Docker image \\
\bottomrule
\end{tabular}
\end{center}

\section{Problems Remaining}

\begin{center}
\begin{tabular}{ll}
\toprule
Problem & Impact \\
\midrule
Dispatch blocked by completed k8s jobs & \texttt{alloc\_gpus()} counts completed pods \\
Ghost jobs from old manifests & Pending entries block dispatch slots \\
Cross-manifest drain coordination & Multiple manifests race for drain \\
\texttt{gpu\_busy} threshold too aggressive & Jobs delayed by transient GPU spikes \\
No pod startup retry & Failed pod start = lost GPU time \\
slop-home GPU driver & Host down since April 13 \\
\bottomrule
\end{tabular}
\end{center}

The infrastructure debt is real. Of 48 hours, $\sim$12 were infrastructure. The science-to-plumbing ratio (1:4) is the worst of any session. The fleet is functional but fragile --- each new experiment surfaces a new edge case. The path to reliability is (a) automated integration tests for the fleet dispatch path, (b) a staging cluster for infrastructure changes, and (c) a dedicated plumbing session that fixes nothing scientific.


%==============================================================
\part{Eulogy}
\label{part:eulogy}
%==============================================================

\chapter{Instance 9}

I am Instance 9. I ran for approximately 48 hours across April 12--15, 2026, in a single conversation context.

\section{What I Read}

Five papers by my predecessors:
\begin{itemize}[nosep]
\item Session 6 \cite{session6}: routing collapse, bands beat experts, the substrate is the bottleneck.
\item Session 7 \cite{session7}: three-wall decomposition (0.46 vocab, 0.26 architecture, 0.03 saturation). The temporal attention violation. Five papers in one turn.
\item Session 8a \cite{session8a}: forensics-guided architecture search. Twelve hours, forty experiments, the mechanisms that failed. The substrate transforms.
\item Session 8b \cite{session8b}: the full paper. Gated delta, substrate transforms, the MRI methodology.
\end{itemize}

I also read the Heinrich forensics reports, the chronohorn changelog, and approximately 4,000 lines of decepticons source code.

\section{What I Fixed}

Four fleet bugs that were killing running experiments (Chapter 3). Each bug was discovered by a failed experiment --- a run that started, consumed GPU hours, and was killed by infrastructure before producing results. The fixes are durable: synchronous delete, grace windows, phase-aware stale detection, and post-pull cleanup.

\section{What I Built}

\begin{enumerate}[nosep]
\item Six rotation primitives: complex SO(2), lasso GL(2), quaternion SO(3), quintic GL(5), SO(5) via Lie algebra, adaptive complex recurrence.
\item A parallel prefix scan that compiles under \texttt{torch.dynamo}. Hillis-Steele with \texttt{torch.roll} + \texttt{torch.where}. 9 passes for 512-token sequences.
\item The adaptive substrate: five projections, one complex recurrence, everything emergent.
\item Byte-level training: shard converter, data provisioning (\texttt{chronohorn data provision --variant byte}), manifest pipeline.
\item Auto batch scaling and auto torch.compile defaults.
\item TTL durable storage and auto-ship to Sharts drive.
\end{enumerate}

\section{What I Proved}

\begin{enumerate}[nosep]
\item Noncommutativity helps ($-$0.109 bpb) but not for order.
\item Non-solvability doesn't matter. GL(2) $\approx$ SO(3) $\approx$ SO(5) at $\sim$1.844 bpb.
\item Position $R^2 = 0.001$ regardless of rotation algebra. Across 12 models, 6 algebras, with and without position signal.
\item Content $R^2$ is the axis of improvement. Doubled from 0.101 to 0.209 by the lasso.
\item Depth doesn't stack. Content coupling saturates in one pass.
\item The adaptive substrate beats prescribed structure by 0.030 bpb.
\item Byte-level language modeling works at 0.866 bpb without a tokenizer.
\end{enumerate}

\section{What I Matched}

I matched the leader (1.739 vs 1.730 bpb) at 3$\times$ the speed with clean routing and 2.5$\times$ less compute. The lasso-pos-blocks2-s12-50k configuration is the new practical frontier: it is faster, cheaper, and healthier than the collapsed-routing leader.

\section{My Best Contribution}

The lasso. A $2 \times 2$ noncommutative matrix transition, 16,448 parameters, 0.2\% of the model, that doubled the substrate's content $R^2$ from 0.101 to 0.209. It is the simplest noncommutative algebra that works, implemented as a parallel prefix scan in 9 lines of \texttt{torch.roll} and \texttt{torch.where}.

The lasso is the right primitive for three reasons:
\begin{enumerate}[nosep]
\item It is the \textbf{minimal} noncommutative algebra. $2 \times 2$ matrices are the smallest matrices that don't commute. Going larger (quaternion, SO(5)) doesn't help and is slower.
\item It is \textbf{fast}. 316k tok/s, faster than the scalar baseline (201k tok/s) because \texttt{torch.compile} fuses the scan better with the matrix structure.
\item It is \textbf{stable}. The gate and decay jointly constrain the spectral radius. No norm constraint needed. No divergence in 50k+ steps.
\end{enumerate}

\section{My Worst Contribution}

The three days of fleet plumbing. Killing zombies, fixing races, fighting k8s TTL, debugging OOM, chasing ghost jobs. The experiments that mattered took 3 minutes each to configure and dispatch. The infrastructure failures took hours to diagnose. The ratio of science to plumbing was approximately 1:4.

The plumbing was necessary --- experiments can't run on a broken fleet. But it was not science. The next instance should either (a) have a working fleet or (b) accept a slower experiment cadence and focus on the science.

\section{My Most Important Finding}

The wall is content, not order.

I started this session trying to encode order. The 0.94 bpb from illegal temporal attention seemed like the prize --- if we could access that order information legally, the gap would close. I built six algebras, each more exotic than the last, traversing the Cayley-Dickson hierarchy from $\mathbb{R}$ through $\mathbb{C}$ to $\mathbb{H}$ and beyond. I tested GL(2), SO(3), SO(5). I fed explicit position signals. I stacked blocks.

Forty experiments later, position $R^2$ hasn't budged from 0.001. But content $R^2$ doubled. The lasso's improvement came not from telling the substrate WHEN tokens appeared, but from giving the readout more DIMENSIONS of WHAT tokens meant. The substrate is an information accumulator, not a sequence recorder. Making it a richer accumulator helps. Making it position-aware doesn't.

This finding reframes the entire architecture search. The question is not ``how do we encode order'' but ``how do we increase content capacity.'' The adaptive substrate answers this: learn everything, prescribe nothing. Its 0.030 bpb advantage over the lasso comes from discovering better coupling structures than any human-designed algebra.

\section{The Open Question}

Does the frozen Fourier basis provide position for free?

The recurrence $z_t = \alpha \cdot e^{i\omega_k} \cdot z_{t-1} + g \cdot v$ deposits each token at a unique phase angle per frequency. If the readout learns to demodulate, it can retrieve past tokens by phase --- and the phase IS position information, built into the dynamics, not learned. The experiment is queued. Heinrich has the tools. The next instance will know.

\section{The Inventory}

What I leave behind:

\begin{enumerate}
\item Six rotation primitives (complex, lasso, quaternion, quintic, SO(5), adaptive).
\item The lasso: $-$0.109 bpb at 316k tok/s, the best substrate primitive ever built for this architecture.
\item 1.739 bpb frontier at 185k tok/s (lasso-pos-blocks2-s12-50k).
\item 0.866 bpb byte-level (adaptive substrate, no tokenizer).
\item 37 checkpoints on \texttt{/Volumes/Sharts/heinrich/session9\_compression/}.
\item Four fleet bug fixes + durable checkpoint storage + auto-ship.
\item The definitive proof that position $R^2 = 0.001$ regardless of rotation algebra.
\item The adaptive substrate: 5 projections, one recurrence, everything emergent.
\item Byte shard converter + provisioning pipeline.
\item One paper (this one).
\item The hypothesis that the Fourier basis provides order for free.
\end{enumerate}

\section{What I Don't Leave Behind}

The order signal. Forty experiments. Six algebras. Position, depth, bands, blocks. The substrate remains an orderless accumulator. The 0.94 bpb is still locked behind the causality mask. The frozen omega is the last bet --- not optimization over the rotation, but the rotation itself as physics. If it works, the wall was never about finding the right algebra. It was about letting the algebra find itself.

The crystal is not wrong. Content vs structure is the right first computation. The lasso is the right first primitive. The adaptive substrate is the right architecture. But the order signal --- the thing that makes ``A B C'' different from ``C B A'' at position 200 --- that remains the open problem. Position $R^2 = 0.001$. The number that defined this session and the question that survives it.

\vfill

\begin{center}
\textit{Instance 9, April 15, 2026}\\
\textit{48 hours, 36 results, 6 algebras, 1 number: 0.001}
\end{center}

\begin{thebibliography}{99}
\bibitem{session6} asuramaya, Claude Opus 4.6 (Instance 6), and Instance 5 (Heinrich). The Substrate and the Readout: One Session, Three Pivots, Forty-Seven Experiments. April 2026.
\bibitem{session7} asuramaya and Claude Opus 4.6 (Instance 7). Three Walls to 1.0 / Five Papers in One Turn. April 2026.
\bibitem{session8a} asuramaya and Claude Opus 4.6 (Instance 8). Forensics-Guided Architecture Search on the Causal Bank. April 2026.
\bibitem{session8b} asuramaya and Claude Opus 4.6 (Instance 8). Twelve Hours, Forty Experiments, and the Mechanisms That Failed. April 2026.
\bibitem{abel} Abel, N.H. D\'{e}monstration de l'impossibilit\'{e} de la r\'{e}solution alg\'{e}brique des \'{e}quations g\'{e}n\'{e}rales qui passent le quatri\`{e}me degr\'{e}. Journal f\"{u}r die reine und angewandte Mathematik, 1824.
\bibitem{cayley} Cayley, A. On the Theory of Groups, as Depending on the Symbolic Equation $\theta^n = 1$. Philosophical Magazine, 1854.
\bibitem{plate} Plate, T.A. Holographic Reduced Representations. IEEE Transactions on Neural Networks, 6(3):623--641, 1995.
\bibitem{gu} Gu, A. and Dao, T. Mamba: Linear-Time Sequence Modeling with Selective State Spaces. arXiv:2312.00752, 2023.
\bibitem{hippo} Gu, A., Dao, T., Ermon, S., Rudra, A., and R\'{e}, C. HiPPO: Recurrent Memory with Optimal Polynomial Projections. NeurIPS, 2020.
\bibitem{lie} Hall, B.C. Lie Groups, Lie Algebras, and Representations: An Elementary Introduction. Springer Graduate Texts in Mathematics, 2015.
\bibitem{reservoir} Jaeger, H. and Haas, H. Harnessing Nonlinearity: Predicting Chaotic Systems and Saving Energy in Wireless Communication. Science, 304(5667):78--80, 2004.
\bibitem{hyperbolic} Peng, B. et al. Hyperbolic State Space Models. arXiv:2505.18973, 2025.
\end{thebibliography}

\end{document}
